\scanpage{53}
\mspage{6}
polygone régulier. Je veux décrire des groupoïdes équivalents aux groupoïdes fondamentaux
\[
(28)\qquad\Pi_1P\Sigma^*\simeq\Pi_1V^{**},\qquad\Pi_1\Sigma^*
\]
par générateurs et relations, avec des sommets bien choisis dans les espaces envisagés, stables par l'action de $D$ ainsi que de celle de $\tau$. Bien sûr, il faut garder présent le morphisme de groupoïdes
\diagram{\begin{tikzcd}[ampersand replacement=\&]
\text{(29)}\quad\Pi_1P\Sigma^*=\Pi_1V^{**}\arrow[r,no head]\&\Pi_1\Sigma^*
\end{tikzcd}}
Sur la situation (28), (29) opère le groupe
\[
\deleted{\unsure{\mathbb C}}\quad D\times\mathbb C^*\times\{1,\tau\}
\]
\margin{$D$ opère tautologiquement $\mathbb C^*$ par homothéties (mais c'est l'action \uline{opposée} de celle dans $P\Sigma^*$)}
\margin{\uline{NB} $\sigma$ n'opère sur $\Sigma^*,P\Sigma^*$ que si $n$ est pair (donc \inserted{pas dans le cas $n=3$} le plus \deleted{\unsure{intéressant}}).}
\margin{\uline{NB} action des groupes \uline{compatibles}, l'élément $(\unsure{\sigma_0},-1)$ opère trivialement}
on s'intéresse à l'action du sous-groupe fini
\[
(30)\qquad D\times\mu_m\times\{1,\tau\}
\]
\margin{d'ordre \deleted{\unsure{4}} $4m^2=16n^2$ ($144$ si $n=3$) mais \unsure{qui en fait a le quotient} fidèle est d'ordre $8n^2$ ($72$ si $n=3$)}
\margin{\uline{NB}. c'est le quotient $(D\times\mu_m)/\{\pm1\}\times\{1,\tau\}$ qui opère fidèlement.}
On peut \unsure{d'ailleurs} prendre le sous-groupe
\[
(31)\qquad D^+\times\mu_m\simeq\mu_m(\underline\omega)\times\mu_m
\]
de (30), en notant que pour la représentation
\mspage{52}
